% Part:sets-functions-relations
% Chapter: sets
% Section: schroder-bernstein

\documentclass[../../../include/open-logic-section]{subfiles}

\begin{document}

\olfileid{sfr}{siz}{sb}

\olsection{आकारमानाची कल्पना आणि श्रेडर--बर्नस्टाइन प्रमेय}

\begin{explain}
पुढील विचार अंतर्ज्ञानाला पटणारा आहे: $A$ हा $B$ पेक्षा आकारमानाने मोठा नाही
आणि $B$ हा $A$ पेक्षा आकारमानाने मोठा नाही, तर $A$ आणि $B$ तुल्यबल
आहेत. खरे सांगायचे तर, हा विचार \emph{चुकीचा} असता, तर तुल्यबलतेच्या
आपल्या व्याख्येचा संचांच्या ``आकारमानां''च्या तुलनेशी काही संबंध आहे,
असे म्हणण्याला आपल्याजवळ जवळजवळ आधारच उरला नसता!{} सुदैवाने हा
अंतःप्रज्ञ विचार बरोबर आहे. श्रेडर--बर्नस्टाइन प्रमेय त्याचे समर्थन
करते.
\end{explain}

\begin{thm}[श्रेडर--बर्नस्टाइन]
	\ollabel{thm:schroder-bernstein}
	$\cardle{A}{B}$ आणि $\cardle{B}{A}$ असेल,
	तर $\cardeq{A}{B}$.
\end{thm}

\begin{explain}
दुसऱ्या शब्दांत, $A$ पासून~$B$ कडे जाणारे !!a{injection} आणि $B$
पासून~$A$ कडे जाणारे !!a{injection} असेल, तर $A$ पासून~$B$ कडे
जाणारे !!a{bijection} असते.

पण या निष्कर्षाची सिद्धता खरोखरच बरीच \emph{कठीण} आहे. कँटर यांनी
निष्कर्ष मांडला असला, तरी इतरांनी तो सिद्ध केला.
\footnote{इतिहासाच्या अधिक माहितीसाठी, उदाहरणार्थ,
\citet[pp.~165--6]{Potter2004} पाहा.}
\oliflabeldef{sfr:cardinals:card-sb:sec}{आपण
श्रेडर--बर्नस्टाइनची \emph{सिद्धता} फक्त
\olref[sfr][cardinals][card-sb]{sec} मध्ये देऊ शकू.}{}%
आत्तासाठी हा निष्कर्ष विश्वासाने स्वीकारावा लागेल.

सुदैवाने श्रेडर--बर्नस्टाइन प्रमेय \emph{बरोबर} आहे आणि म्हणून आपण
परिभाषित केलेले $\cardeq{A}{B}$ आणि $\cardle{A}{B}$ हे संबंध
``आकारमाना''शी निगडित आहेत, असे मानणे समर्थित होते. शिवाय
श्रेडर--बर्नस्टाइन प्रमेय फार \emph{उपयुक्त} आहे. दोन तुल्यबल संचांमध्ये
!!a{bijection} शोधणे कठीण असू शकते. या प्रमेयामुळे तुलना दोन दिशांत
विभागता येते: दोन्ही दिशांतील !!a{injection} स्वतंत्रपणे शोधावी
लागतात.
\end{explain}
\end{document}
